% ======================================================================
% Section 3 — The type-mismatch mechanism and four instances
% ======================================================================

\section{The type-mismatch mechanism and four instances}\label{sec:tm}

\subsection{The mechanism}\label{sec:tm-mech}

Call a statement $R$ a \emph{lossless reformulation} of $\lrc_n$ if $R$
holds for all speed vectors at rung $n$ if and only if $\lrc_n$ does, with
no slack in either direction. Call a statement $S$ a \emph{toolkit
strengthening} of $R$ if (i) $S$ implies $R$, strictly, and (ii) the
strictness is of the kind that a named toolkit can exploit: a uniform
margin over a trivial baseline, a spectral gap, an interior rate, a
monotone law in the scale parameter, a saturated bound with slack. The
distinction is not aesthetic; it is the difference between statements
that classification, spectral, harmonic, or polyhedral instruments can
consume and statements they cannot.

Three observations, jointly verified four times in this program, make up
the mechanism.

\begin{enumerate}[leftmargin=1.6em,itemsep=2pt]
\item \textbf{Losslessness transfers truth exactly.} A proof of the
reformulated statement is ipso facto a proof of $\lrc_n$, and conversely.
A rewrite may still change tractability---the program's own finitization
is an equivalent rewrite that made the problem computable instance by
instance---but it cannot discard the conjecture's content: no lossless
rewrite can be simultaneously strictly stronger than $\lrc_n$ and true
at the extremal instances.
\item \textbf{The strengthening is false where it matters.} The
toolkit-relevant strengthenings are strictly stronger than $\lrc$ (they
imply it; the converse fails), and at the extremal instances (the
harmonic and rung-2 families, the reflection-symmetric cascade families,
the tight covering configurations) the extra margin they demand does not
exist. When the demand is a definite statement, it is false, and small
$n$ exhibits the falsity with exact witnesses.
\item \textbf{The instruments land on the closed side.} The instruments a
lossless reformulation naturally supports---exact counts,
classifications, invariances, exhaustiveness checks---compose into
statements that close everything \emph{except} the missing margin: by
theorem in Instance 2 (Lemma~\ref{lem:monomial}), by audit in Instance 1
(the inventory classification), and by measurement in Instances 3 and 4
(the saturation and decay records below).
\end{enumerate}

The four subsections that follow develop one instance each. Three of
them sit on the same pair-sum cascade (an inventory audit, an instrument
closure, a rigidity falsification) and the fourth is independent of it
(the zonotope); what is common to all four is the mechanism, and we
state it once, here, as the paper's organizing claim: \emph{in each
setting, the falsity or unavailability of the strengthening is not an
accident of the example but a consequence of the losslessness that made
the reformulation faithful.}

\subsection{Instance 1: the finite covering reformulation}\label{sec:tm-covering}

The pair-sum equality theorem (\S\ref{sec:positive}) reduces
$M(V)$ to a finite maximum over cyclic groups: $M(V)=\max_{\{u,w\}}
\tau(u,w)$, where $\tau$ is the optimal loneliness on the pair-sum
lattice of denominator $N=v_u+v_w$. Under this reduction $\lrc_n$
becomes the statement that for every $V$ some pair's bad sets
$B_p=\{k\in\Z_N:(n{+}1)\wn{v_pk}<N\}$, $p\notin\{u,w\}$, fail to
cover $\Z_N$: a finite covering problem, exact, and apparently made for
the covering toolkit.

The toolkit delivered. The program proved an exact size law for the bad
sets, a polygonal witness lemma, moat-counting and coincidence-counting
lemmas, a lift law, and a class-group reduction at prime moduli; it
classified the unit covering configurations completely at the rigid
primes (Theorem~\ref{thm:rigid}); and on the full corpus of
$3{,}712{,}576$ primitive $5$-sets with speeds $\le 56$ it closed
$75.6\%$ of instances by proved lemmas alone ($78.2\%$ with the
coincidence lemma) and all of them by per-modulus verification,
localizing the residual obstruction to a single named class
\cite{prog-artifacts}. What it did not deliver, at any point in the
program's multi-year run, is a single statement of the type the residual
obstruction requires: a \emph{uniform growth-in-$k$ margin over a trivial
or measured baseline}.

This is a verifiable claim, and it was verified by audit: the full proof
inventory of the program---$34$ named results across the pair-sum,
group-ring, tensor, and cascade sections---was classified
statement-by-statement. Every proved statement is a closure, an
invariance, an exhaustiveness, or a machine verification. The one
candidate that reads like an exception, the cut law's
``$k$-cancellation,'' resolves as a $k$-scale \emph{invariance} of the
composed bound: it cancels exactly the $k$-dependence that a margin
would need to grow---the opposite of the missing type. The
spacing/packing lemma, the chain-class volume law, and the dilate theorem
were each examined as potential margin carriers and each resolves the
same way: a $k$-growing \emph{count} whose margin over its own trivial
baseline is $k$-free, a degrading order-factor slack, or a conditional
consequence of the hypothesis stack it was meant to prove.

The empirical half of the reading (``no proved margin-type statement
exists in the inventory'') is thereby settled. The structural half (``no
classification toolkit \emph{can} produce one'') is a hypothesis about
tool reach; it is what Instances 2--4 test, in three independent
settings, and what the falsity archive of \S\ref{sec:witness} completes
on the metric side.

\subsection{Instance 2: the transfer-matrix route, closed by theorem}
\label{sec:tm-transfer}

The one margin-shaped candidate instrument in the entire program was a
transfer-matrix (automaton) model of the solution cascade: transfer
operators naturally emit spectral quantities uniform in the scale
parameter, and a uniform spectral gap would be exactly the missing type.
The session that built it fixed its success criterion before computing:
the whole-cascade survival fraction must be bounded by
$\rho^{\mathrm{depth}}$ times the initial mass, with spectral radius
$\rho\le 1-\epsilon$ and $\epsilon$ independent of scale.

\begin{lemma}[Monomial dichotomy; Lemma M of the program records]
\label{lem:monomial}
Every constrained transfer matrix of the lossless cascade is a partial
permutation matrix. Its spectrum is $\{0\}$ together with roots of unity;
its spectral radius is exactly $0$ or $1$, never interior, at any scale,
in any state formulation. The affine pair maps compose to the identity
around any closed fiber loop (the dilation ratios telescope), so the
cyclic transfer product is diagonal, and its trace is the cascade's own
survivor count.
\end{lemma}

The proof is short once the state space is fixed: the cascade is
non-consuming (pure intersection semantics), the drift maps are
bijections, and a bijection constrained to an invariant subset is a
partial permutation. The machine record is not short: $1500/1500$ random
constrained systems; $12{,}000$ exact trace fingerprints over $300$
systems; $66$ homogenized pair systems on the committed families, with
spectral radius $1$ at exactly the six reflection-symmetric families
(their deletion sets close into $2$-cycles and the two-fiber survival
fraction is exactly $1.0000$ at all depths) and $0$ at the other sixty
(total extinction within three iterates, the intermediates set by lattice
runs of solution points, not by geometric decay). The criterion is
unattainable by theorem, not merely unattained.

Two further doors were checked. Coarse-grained or weighted automata can
carry interior spectral radii numerically, but the radius dominates every
cycle's geometric mean, so a uniform gap would force uniform decay of the
census weights---the margin as an \emph{input}, not an output---and the
converse fails (a $2\times2$ witness with all cycle means $\le0.9$ and
radius $1.122$: branching resurrects). Averaged or randomized operators
collapse to the independent-extension baseline. And the reduction closes
the loop: the whole-cascade survival fraction equals the product of
consecutive conditional ratios, each bounded by the pair-volume and
solution-set thinness hypotheses the route was commissioned to prove. The
instrument is the census in different notation; it consumes the
hypotheses it was built to generate. Losslessness forbids generated
decay---bijections preserve measure, deletion filters without
dissipating---and the reformulation is bounded by its own fidelity.

\subsection{Instance 3: character rigidity and saturated envelopes}
\label{sec:tm-char}

The harmonic-analytic route phrased the cascade obstruction in terms of
\emph{characters}: the four singleton coordinates and six pairwise
differences of the placement vector form a closed system of ten
functionals (Lemma C10 of the program records), closed under the
renormalization action up to affine shifts, whose aggregate
marginal-multiset is an orbit invariant. Any class-level census law must
live on this system; the route's natural strengthenings are laws about
its per-character profiles.

They are false, exactly. The singleton-profile orbit invariance (Lemma
V$'$) fails at both fully-enumerated primes: per-key counts match across
orbit members only after a size permutation, and the profiles differ. The
pairwise-profile invariance (V$''$) fails with $0/32$ sigma-matched
tuples. The mechanism is \emph{anchor mixing}: the renormalization
bijection couples every moving slot to the anchor coordinate; slot
singletons map to anchor pairs and pairwise differences to non-anchor
singletons, so no per-character statistic can be invariant---only the
aggregate is.

Worse, the per-character envelopes are \emph{saturated} exactly where the
margin would be needed. On the committed top orbits the pinning envelope
equals the modulus at both $k=2$ primes and $27/29$ at $k=3$ (fullness
fractions $1.000$, $1.000$, $0.931$): the marginal boxes are full, the
$\gamma$-envelope sits at or above the threshold value $4$ at $k=2$, and
the only invariant sparsity is \emph{joint} (median joint densities in
the marginal box: $1.5\%$, $7.4\%$, $6.9\%$). A bound phrased on
individual characters is tight---zero slack---precisely at the
configurations the obligations target.

Finally, the trivial bound sits exactly at the cascade threshold. The
baseline $|Sol|\le p^{m-1}$ realizes the constant
$c_1=(T/(T-6)+o(1))^{m-1}$ (here and below $m=T-3$ is the arc count of
the frontier cell, so $m-1=T-4$)---numerically $256$, $243$, $244.14$ at
$T=8,9,10$---which is exactly the value at which the depth denominator
of the conditional scaling theorem vanishes. The hypothesis H1g's
content is to beat the trivial bound by a factor growing in $k$; the
measured factors ($95\times$, $1{,}207\times$, $3{,}803\times$ at
$k=2,3,4$) grow, but the first of them starts \emph{below} the
threshold constant ($95<244$), which is precisely why the margin must
grow in $k$ and why $k=2$ alone cannot certify anything; the proved
factors do not exist, and the sharp-constant question and the margin
question are one criterion stated twice. This is the type mismatch in its purest
form: the character system is exactly rich enough to locate the
obstruction and exactly too poor to bound it.

\subsection{Instance 4: the Lonely Runner zonotope}\label{sec:tm-zono}

The fourth instance is metric, and its falsity archive is the paper's
sharpest. The quotient-torus model of \S\ref{sec:prelim} gives two
statements: the $\delta$-form (the center class is within
$\mathrm{thr}_n$ of the lattice), which is equivalent to $\lrc_n$, and
the $\rho$-form (every class is), which is strictly stronger and would
open the geometry-of-numbers toolkit---transference bounds, covering
radius inequalities, dilution arguments. The natural lemma a
covering-radius toolkit wants first is exactly $\rho\le\mathrm{thr}_n$.

It is false, with exact witnesses, from $n=5$. At the extremal harmonic
instance $v=(1,2,3,4,5)$, the rational point
$c=(\tfrac1{32},\tfrac{15}{32},\tfrac3{32},\tfrac78)$ has
$D(c)=97/288>\tfrac13=\mathrm{thr}_5$ (the branch-and-bound follow-up
found $607/1792$ at
$c=(\tfrac{11}{256},\tfrac{163}{256},\tfrac{15}{256},
\tfrac{79}{256})$, and the audit of this paper found
$4183/12288$ at
$c=(\tfrac{465}{4096},\tfrac{605}{1024},\tfrac{1019}{4096},
\tfrac{315}{4096})$, and the third certification session closed the row
exactly: $\rho(1,2,3,4,5)=16/47$ at the balance-family witness
$c=(\tfrac{46}{47},\tfrac{92}{235},\tfrac{231}{235},\tfrac{35}{47})$,
\S\ref{sec:bb}): the deepest hole is not the center, and the
covering-radius statement fails at the very instance where the
$\delta$-form is tight---with the failure now measured exactly. At $n=6$ the harmonic witness
$c=(0,\tfrac5{16},\tfrac{15}{16},\tfrac34,\tfrac12)$ gives
$D=29/80>\tfrac5{14}=\mathrm{thr}_6$, and the non-extremal instance
$v=(1,2,3,4,5,7)$---whose center depth is the \emph{computed}
$\delta=1/3$, strictly below $\mathrm{thr}_6$ with slack---carries
$c=(\tfrac1{16},\tfrac9{16},0,\tfrac38,0)$
with $D=4/11>\mathrm{thr}_6$: the falsity is not an extremal-curiosity
confined to the tight family, and the covering-radius failure there is
not an artifact of center-class tightness. The rung-2 instance
$v=(1,2,3,4,5,12)$ carries $c=(\tfrac1{16},\tfrac{11}{16},\tfrac18,
\tfrac7{16},\tfrac78)$ at exactly $D=5/14=\mathrm{thr}_6$: a
zero-slack boundary case that already separates the deepest hole from
the center ($\delta=9/26<5/14$) without strictly exceeding the
threshold. All witnesses are non-torsion points---the deepest holes are
not the center classes---and all are single exact evaluations of
\eqref{eq:runnerform}, checkable by hand (Appendix~\ref{sec:self}).

For $n\le4$ the $\rho$-form holds on every instance tested: at $n=3$
the deepest-hole property (that the maximum of $D$ is attained at the
center class) holds exactly when $3\mid m$ in the family
$(1,2,m)$ for $m\le24$, and at $n=4$ it holds when $4\mid m$ for
$m\le16$ (away-certified at $m=4,8,12$; branch-and-bound certified at
$m=16$, \S\ref{sec:bb}), with exact witnesses against every other $m$.
The two forms have the same truth value throughout the tested small
range---which is precisely why the stronger form was tempting; the
$n=5$ witness is precisely why it is wrong.

The companion hope---that the zonotope's Ehrhart theory could control
the covering radius through the positivity structure of its
$h^*$-polynomial---fails independently, and its failure mode is the
cleanest illustration of type mismatch in the whole program. The
The $h^*$-polynomials are real-rooted, log-concave, and Newton, uniformly,
in every instance tested through $n=10$ in both core families
(Appendix~\ref{sec:hstar}). But the log-concavity \emph{strictness} $S$
decays roughly like $30/n$ in the measured range ($nS$ drifting from
$\approx37$ to $\approx30$ in the harmonic family over
$n=3,\dots,10$, and from $60.5$ to $\approx30$ in the rung-2 family),
and Newton's inequalities force $S$ above a floor $\approx1+4/d$ that
itself declines, so much of the decay is generic and the $1/n$ behavior
cannot persist indefinitely---there is no growing quantitative margin to
spend; and across the instance sweeps the strictness has no
statistically distinguishable correlation with the covering margins at
the achievable sample sizes (Spearman $+0.32$, $p=0.32$, $N=12$ at
$n=3$; $-0.31$, $p=0.45$, $N=8$ at $n=4$; the sign flip is consistent
with noise; these arrays are too small to establish correlation or its
absence). A
coefficient inequality of a counting polynomial at growing scale $t$ is
the wrong type to pin a fixed-scale metric hole depth; the computation
says so twice, once by decay and once by the absence of a measurable
signal, and the covering-radius statement it was meant to feed is false
at the extremal instances anyway. There is no mechanism, and none is
needed: the bridge is dead on both ends.

\subsection{The pattern}

Table~\ref{tab:instances} assembles the four instances. Each row names
the reformulation, the toolkit strengthening it wanted, the status of
that strengthening, and the witness. The rows are independent; the last
column is always finite, exact, and small. We emphasize what the table
does \emph{not} contain: any statement falsifying $\lrc$ itself, or any
statement weakening our certainty that the surviving equivalent forms
(the $\delta$-form, the pair-sum form) are exactly as hard as the
original.

\begin{table}[htbp]
\centering
\small
\caption{The four instances of type mismatch. ``Margin type'' is the
statement class the toolkit needed; ``status'' is its verified fate,
with the kind of verification distinguished: \emph{false} (an exact
witness falsifies a definite statement), \emph{closed by theorem} (a
theorem shows the instrument cannot produce the output), and
\emph{zero in inventory} (an audit finding: no statement of the type
exists among the proved results). The first three rows sit on the same
pair-sum cascade; the fourth is the independent zonotope setting.}
\label{tab:instances}
\begin{tabularx}{\textwidth}{@{} l >{\raggedright\arraybackslash}X
>{\raggedright\arraybackslash}X l @{}}
\toprule
Reformulation & Toolkit wanted & Status & Witness \\
\midrule
Pair-sum covering & growth-in-$k$ margin & zero in inventory (audit;
$34$ statements); $k$-cancellation is a $k$-invariance & audit record
($34$ statements) \\
Transfer cascade & spectral gap $1-\epsilon$ & closed by theorem
(Lemma~M; survival $=1$ at reflection families) & six families, exact \\
Character rigidity & per-character profile laws & false (V$'$/V$''$
witnesses); envelopes saturated; trivial bound at threshold &
exhaustive orbits; $0/32$ \\
LR-zonotope & $\rho\le\mathrm{thr}_n$ (covering) & false from $n=5$;
$h^*$ strictness decays, no signal & $16/47$ exact; $29/80$, $4/11$,
$17/47$, $455/1269$ \\
\bottomrule
\end{tabularx}
\end{table}
